\chapter*{Supporting Technologies}
\noindent
The following third-party hardware and software were used to build and evaluate CARS. The contribution of this project is the CARS design, its controller engine, and the evaluation; the items below are the platform it was built on.

\vspace{0.5em}
\noindent\textbf{Hardware}
\begin{itemize}
\item Three workstations forming the testbed: one runs the Open vSwitch fabric (\texttt{ovs1} and \texttt{ovsgw}), the Snort sensor and the supporting services, and hosts the supervisory and attacker endpoints; one runs the CARS controller; one runs the second cell (\texttt{ovs2}).
\item Two Siemens S7-1200 programmable logic controllers driving the physical process. PLC1 is a 1212C, model 6ES7~212-1BE40-0XB0, firmware 4.2.3, confirmed by scanning it during evaluation.
\item A Siemens KTP700 operator panel as the human-machine interface on the process cell.
\item A hAP access point providing the isolated wired management network that carries the OpenFlow control plane, kept separate from the operational-technology data plane.
\item A Windows PC acting as the engineering workstation, running TIA Portal and Factory IO.
\end{itemize}

\noindent\textbf{Software}
\begin{itemize}
\item TIA Portal, used to program the PLCs and the OB30 tank-control logic.
\item Factory IO with its S7-1200 driver, providing the hardware-in-the-loop tank-level process wired to the live PLC.
\item Open vSwitch, providing the software-defined data plane (the switches \texttt{ovsgw}, \texttt{ovs1} and \texttt{ovs2}) under OpenFlow~1.3 control.
\item os-ken, the OpenFlow~1.3 controller framework hosting the CARS engine (\texttt{cars\_engine.py}); the engine and the supporting tools are written in Python.
\item Snort, used as the deep-packet-inspection sensor that recovers the industrial operation from S7 and Modbus traffic.
\item python-snap7, the S7 client library used by the process model, the remediation agent and the attack clients.
\item Linux network namespaces, used to isolate the supervisory endpoints (SCADA, Modbus operator, historian, remediation agent) on the fabric-and-sensing node, each bound to a switch port.
\item Kali Linux, used as the external attacker virtual machine attached to the operational segment, together with nmap for reconnaissance and a set of custom Python clients for the S7 and Modbus write, false-data-injection and fragmentation tests.
\item tcpdump, with Wireshark and tshark, used to capture and analyse traffic at the wire points for the packet-level evaluation.
\end{itemize}
